\chapter{Formation of Singularities}
\label{chap:cz-singularity-formation}
\section{Cheeger Type Compactness}
\begin{definition}[pointed smooth convergence of Riemannian manifolds]
  \label{def:cz-pointed-smooth-convergence-of-riemannian-manifolds}
  \source{cao-zhu}
  \dcref{Definition 4.1.1}
  \group{Cao--Zhu Cheeger Type Compactness}
Let $(M_k,g_k,p_k)$ be a sequence of marked complete Riemannian
manifolds, with metrics $g_k$ and marked points $p_k\in M_k$. Let
$B(p_k, s_k)\subset M_k$ denote the geodesic ball centered at
$p_k\in M_k$ and of radius $s_k$ ($0 < s_k \leq +\infty$). We say
a sequence of marked geodesic balls $(B(p_k,s_k),g_k, p_k)$ with
$s_k \rightarrow s_{\infty}(\leq +\infty)$ {\bf converges} in the
$C_{loc}^{\infty}$ topology {\bf to a marked} (maybe noncomplete)
{\bf manifold}\index{converges to a marked manifold}
$(B_{\infty},g_{\infty},p_{\infty})$, which is an open geodesic
ball centered at $p_{\infty}\in B_{\infty}$ and of radius
$s_{\infty}$ with respect to the metric $g_{\infty}$, if we can
find a sequence of exhausting open sets $U_k$ in $B_{\infty}$
containing $p_{\infty}$ and a sequence of diffeomorphisms $f_k$ of
the sets $U_k$ in $B_{\infty}$ to open sets $V_k$ in $B(p_k,s_k)
\subset M_k$ mapping $p_{\infty}$ to $p_k$ such that the pull-back
metrics $\tilde{g}_k=(f_k)^*g_k$ converge in $C^\infty$ topology
to $g_{\infty}$ on every compact subset of $B_{\infty}$.
\end{definition}
\begin{theorem}[Cheeger-type pointed compactness]
  \label{thm:cz-cheeger-type-pointed-compactness}
  \source{cao-zhu}
  \uses{cor:cz-injectivity-radius-from-local-curvature-scale, prop:cz-local-metric-subsequence-convergence, prop:cz-transition-isometry-subsequence-convergence}
  \dcref{Theorem 4.1.2}
  \group{Cao--Zhu Cheeger Type Compactness}
Let\; $(M_k,\,g_k,\,p_k)$\; be\; a\; sequence\; of marked complete
Riemannian manifolds of dimension $n$. Consider a sequence of
geodesic balls $B(p_k,s_k) \subset M_k$ of radius $s_k$ $(0<s_k\le
\infty)$, with $s_k \rightarrow s_{\infty}(\le \infty)$, around
the base point $p_k$ of $M_k$ in the metric $g_k$. Suppose
\begin{itemize}
\item[(a)] for every radius $r<s_{\infty}$ and every integer $l\ge
0$ there exists a constant $B_{l,r}$, independent of $k$, and
positive integer $k(r,l) < +\infty$ such that as $k \geq k(r,l)$,
the curvature tensors $Rm(g_k)$ of the metrics $g_k$ and their
$l^{\mbox{th}}$-covariant derivatives satisfy the bounds
$$
|\nabla^lRm(g_k)|\leq B_{l,r}
$$
on the balls $B(p_k, r)$ of radius $r$ around $p_k$ in the metrics
$g_k$; and \item[(b)] there exists a constant $\delta>0$
independent of $k$ such that the injectivity radii
inj$\,(M_k,p_k)$ of $M_k$ at $p_k$ in the metric $g_k$ satisfy the
bound
$$
{\rm inj}\,(M_k,p_k)\geq\delta.$$
\end{itemize}
Then there exists a subsequence of the marked geodesic balls
$(B(p_k,\,s_k),\,g_k,$ $p_k)$ which converges to a marked geodesic
ball $(B(p_{\infty},s_{\infty}),g_{\infty},p_{\infty})$ in
$C^\infty_{\rm loc}$ topology. Moreover the limit is complete if
$s_{\infty} = +\infty$.

\end{theorem}
\begin{proposition}[local metric subsequence convergence]
  \label{prop:cz-local-metric-subsequence-convergence}
  \source{cao-zhu}
  \dcref{Claim 1}
  \group{Cao--Zhu Cheeger Type Compactness}
By passing to another subsequence we
can guarantee that for each $\alpha$ (and indeed all $\alpha$ by
diagonalization) the metrics $g_k^{\alpha}$ (or $\bar
g_k^{\alpha}$ or ${\bar{\bar g}}_k^{\alpha}$) converge uniformly
with their derivatives to a smooth metric $g^{\alpha}$ (or $\bar
g^{\alpha}$ or ${\bar{\bar g}}^{\alpha}$) on $E^{\alpha}$ (or
$\bar E^{\alpha}$ or ${\bar{\bar E}}^{\alpha}$) which is also in
geodesic coordinates.
\end{proposition}
\begin{proposition}[transition isometry subsequence convergence]
  \label{prop:cz-transition-isometry-subsequence-convergence}
  \source{cao-zhu}
  \dcref{Claim 2}
  \group{Cao--Zhu Cheeger Type Compactness}
The isometries $J_k^{\alpha\beta}$
(and $\bar J_k^{\alpha\beta}$ and $J_k^{\beta\alpha}$ and $\bar
J_k^{\beta\alpha}$) always have a convergent subsequence. \vskip
0.2cm
\end{proposition}
\begin{definition}[pointed convergence of evolving manifolds]
  \label{def:cz-pointed-convergence-of-evolving-manifolds}
  \source{cao-zhu}
  \dcref{Definition 4.1.3}
  \group{Cao--Zhu Cheeger Type Compactness}
Let $(M_k,g_k(t),p_k)$ be a sequence of evolving marked complete
Riemannian manifolds, with the evolving metrics $g_k(t)$ over a
fixed time interval $t\in(A,\Omega]$, $A < 0 \leq \Omega$, and
with the marked points $p_k\in M_k$. We say a sequence of evolving
marked $(B_0(p_k,s_k),g_k(t),p_k)$ over $t \in (A,\Omega]$, where
$B_0(p_k,s_k)$ are geodesic balls of $(M_k,g_k(0))$ centered at
$p_k$ with the radii $s_k \rightarrow s_{\infty}(\leq +\infty)$,
{\bf converges} in the $C_{\rm loc}^{\infty}$ topology {\bf to an
evolving marked} (maybe noncomplete) {\bf
manifold}\index{converges to an evolving marked manifold}
$(B_{\infty},g_{\infty}(t),p_{\infty})$ over $t\in(A,\Omega]$,
where, at the time $t=0$, $B_{\infty}$ is a geodesic open ball
centered at $p_{\infty}\in B_{\infty}$ with the radius
$s_{\infty}$, if we can find a sequence of exhausting open sets
$U_k$ in $B_{\infty}$ containing $p_{\infty}$ and a sequence of
diffeomorphisms $f_k$ of the sets $U_k$ in $B_{\infty}$ to open
sets $V_k$ in $B(p_k,s_k) \subset M_k$ mapping $p_{\infty}$ to
$p_k$ such that the pull-back metrics
$\tilde{g}_k(t)=(f_k)^*g_k(t)$ converge in $C^\infty$ topology to
$g_{\infty}(t)$ on every compact subset of
$B_{\infty}\times(A,\Omega]$.
\end{definition}
\begin{lemma}[time-uniform smooth convergence from one time slice]
  \label{lem:cz-time-uniform-smooth-convergence-from-one-time-slice}
  \source{cao-zhu}
  \dcref{Lemma 4.1.4}
  \group{Cao--Zhu Cheeger Type Compactness}
Let $(M,g)$ be a Riemannian manifold, $K$ a compact subset of $M$,
and $\tilde{g}_k(t)$ a collection of solutions to Ricci flow
defined on neighborhoods of $K\times[\alpha,\beta]$ with
$[\alpha,\beta]$ containing $0$.  Suppose that for each $l\geq0$,
\begin{itemize}
\item[(a)] $C_0^{-1}g\leq\tilde{g}_k(0)\leq C_0g,\ \ on\ K,\ for\
all\ k,$ \item[(b)] $|\nabla^l\tilde{g}_k(0)|\leq C_l,\ \ on\ K,\
for\ all\ k,$ \item[(c)]
$|\tilde{\nabla}^l_kRm(\tilde{g}_k)|_k\leq C'_l,\ on\
K\times[\alpha,\beta],\ for\ all\ k,$
\end{itemize}
for some positive constants $C_l,\ C'_l,\ l=0,1,\ldots,$
independent of $k$, where $Rm(\tilde{g}_k)$ are the curvature
tensors of the metrics $\tilde{g}_k(t)$, $\tilde{\nabla}_k$ denote
covariant derivative with respect to $\tilde{g}_k(t)$, $|\cdot|_k$
are the length of a tensor with respect to $\tilde{g}_k(t)$, and
$|\cdot|$ is the length with respect to $g$. Then the metrics
$\tilde{g}_k(t)$ satisfy
$$
\tilde{C_0}^{-1}g\leq\tilde{g}_k(t) \leq \tilde{C_0}g,\ \ on\
K\times[\alpha,\beta]
$$
and
$$
|\nabla^l\tilde{g}_k|\leq\tilde{C_l},\ on\ K\times[\alpha,\beta],\
\ l=1,2,\ldots,
$$
for all $k$, where $\tilde{C_l},\ l=0,1,\ldots,$ are positive
constants independent of $k$.
\end{lemma}
\begin{theorem}[Hamilton compactness for Ricci flows]
  \label{thm:cz-hamilton-compactness-for-ricci-flows}
  \source{cao-zhu}
  \dcref{Theorem 4.1.5}
  \group{Cao--Zhu Cheeger Type Compactness}
Let $(M_k,g_k(t),p_k),$ $t\in(A,\Omega]$ with $A<0 \leq \Omega$,
be a sequence of evolving marked complete Riemannian manifolds.
Consider a sequence of geodesic balls $B_0(p_k,s_k) \subset M_k$
of radii $s_k (0<s_k \leq +{\infty})$, with $s_k \rightarrow
s_{\infty}$ $(\leq +{\infty})$, around the base points $p_k$ in
the metrics $g_k(0)$. Suppose each $g_k(t)$ is a solution to the
Ricci flow on $B_0(p_k,s_k) \times (A,\Omega]$. Suppose also
\begin{itemize}
\item[(i)] for every radius $r<s_{\infty}$ there exist positive
constants $C(r)$ and $k(r)$ independent of $k$ such that the
curvature tensors $Rm(g_k)$ of the evolving metrics $g_k(t)$
satisfy the bound
$$
|Rm(g_k)|\leq C(r),
$$
on $B_0(p_k,r) \times (A,\Omega]$ for all $k \geq k(r)$, and
\item[(ii)] there exists a constant $\delta > 0$ such that the
injectivity radii of $M_k$ at $p_k$ in the metric $g_k(0)$ satisfy
the bound
$$
{\rm inj}\,(M_k,p_k,g_k(0))\geq\delta>0,
$$
for all $k = 1, 2, \ldots$.
\end{itemize}
Then there exists a subsequence of evolving marked
$(B_0(p_k,s_k),g_k(t),p_k)$ over $t \in (A,\Omega]$ which converge
in $C^\infty_{loc}$ topology to a solution
$(B_{\infty},g_{\infty}(t),p_{\infty})$ over $t \in (A,\Omega]$ to
the Ricci flow, where, at the time $t=0$, $B_{\infty}$ is a
geodesic open ball centered at $p_{\infty}\in B_{\infty}$ with the
radius $s_{\infty}$. Moreover the limiting solution is complete if
$s_{\infty}=+\infty$.
\end{theorem}
\section{Injectivity Radius Estimates}
\begin{theorem}[Gromoll–Meyer injectivity-radius estimate]
  \label{thm:cz-gromollmeyer-injectivity-radius-estimate}
  \source{cao-zhu}
  \dcref{Theorem 4.2.1}
  \group{Cao--Zhu Injectivity Radius Estimates}
\index{Gromoll-Meyer injectivity radius estimate} Let $M$
be a complete, noncompact Riemannian manifold with positive
sectional curvature. Then the injectivity radius of $M$ satisfies
the following estimate
$$
{\rm inj}\,(M)\geq\frac{\pi}{\sqrt{K_{\max}}}.
$$
\end{theorem}
\begin{theorem}[local injectivity-radius estimate from volume and curvature]
  \label{thm:cz-local-injectivity-radius-estimate-from-volume-and-curvature}
  \source{cao-zhu}
  \dcref{Theorem 4.2.2}
  \group{Cao--Zhu Injectivity Radius Estimates}
Let $B(x_0, 4r_0)$, $0<r_0<\infty$, be a geodesic ball in an
$n$-dimensional complete Riemannian manifold $(M,g)$ such that the
sectional curvature $K$ of the metric $g$ on $B(x_0,4r_0)$
satisfies the bounds
$$
\lambda\leq K\leq\Lambda
$$
for some constants $\lambda$ and $\Lambda$. Then for any positive
constant $r\le r_0$ $($we will also require $r\le
\pi/(4\sqrt{\Lambda})$ if $\Lambda>0)$ the injectivity radius of
$M$ at $x_0$ can be bounded from below by
$$
\mbox{\rm inj}(M, x_0)\geq r\cdot\frac{\Vol(B(x_0,
r))}{\Vol(B(x_0,r))+V^n_{\lambda}(2r)},
$$
where $V_{\lambda}^n(2r)$ denotes the volume of a geodesic ball of
radius $2r$ in the $n$-dimensional simply connected space form
$M_{\lambda}$ with constant sectional curvature $\lambda$.
\end{theorem}
\begin{corollary}[injectivity radius from local curvature scale]
  \label{cor:cz-injectivity-radius-from-local-curvature-scale}
  \source{cao-zhu}
  \dcref{Corollary 4.2.3}
  \group{Cao--Zhu Injectivity Radius Estimates}
Suppose $B(p_0, s_0)$ $(0<s_0\le \infty)$ is a geodesic ball in an
$n$-dimensional complete Riemannian manifold $M$ having the
property that for any $0<s<s_0$ the sectional curvature $K$ on
$B(p_0, s)$ satisfies the bound
$$
|K|\le \Lambda(s)
$$
for\; some\; positive\; increasing\; function\; $\Lambda$\;
defined\; on\; $[0,\,s_0)$.\; Then\; for\; any\; point\; $p\;\in\;
B(p_0,\, s_0)$\; with\; $d(p_0,\, p)\;=\;s$\; and\; any\;
positive\; number $r\le \min\{\pi/4\sqrt{\Lambda(s+2r_0)},r_0\}$
with $r_0=(s_0-s)/4$, the injectivity radius of $M$ at $p$ is
bounded below by
\begin{align*}
& {\rm inj}\,(M,p) \\
&\geq r \frac{V^{n}_{-\Lambda(s+2r_0)}(r)\cdot
\Vol(B(p_0,r_0))}{V^{n}_{-\Lambda(s+2r_0)}(r)\Vol(B(p_0,r_0))+
V_{-\Lambda(s+2r_0)}^n(2r)V^{n}_{-\Lambda(s+2r_0)}(s+2r_0)}.
\end{align*}
In particular, we have \be
\inj(M,p)\geq \rho_{n, \delta, \Lambda}(s) %\eqno(4.2.2)
\ee where $\delta>0$ is a lower bound of the injectivity radius
$\inj(M, p_0)$ at the origin $p_0$ and $\rho_{n,\delta, \Lambda}:
[0, s_0) \to \mathbb{R^+}$ is a positive decreasing function that
depends only on the dimension $n$, the lower bound $\delta$ of the
injectivity radius $\inj(M,p_0)$, and the function $\Lambda$.
%$$\eqno \#$$ \vskip 0.3cm
\end{corollary}
\begin{theorem}[little loop lemma]
  \label{thm:cz-little-loop-lemma}
  \source{cao-zhu}
  \dcref{Theorem 4.2.4}
  \group{Cao--Zhu Injectivity Radius Estimates}
Let $g_{ij}(t)$, $0\le t<T<+\infty$, be a solution of the Ricci
flow on a compact manifold $M$. Then there exists a constant
$\rho>0$ having the following property: if at a point $x_0\in M$
and a time $t_0\in [0,T)$,
$$ |Rm|(\cdot,t_0)\leq r^{-2} \quad
\text{on}\ B_{t_0}(x_0,r)
$$
for some $r\leq \sqrt{T}$, then the injectivity radius of $M$ with
respect to the metric $g_{ij}(t_0)$ at $x_0$ is bounded from below
by
$$
\inj(M,x_0,g_{ij}(t_0))\geq \rho r.
$$
%$$\eqno \#$$
\end{theorem}
\section{Limiting Singularity Models}
\begin{definition}[almost maximum curvature points]
  \label{def:cz-almost-maximum-curvature-points}
  \source{cao-zhu}
  \dcref{Definition 4.3.1}
  \group{Cao--Zhu Limiting Singularity Models}
We say that $\{(x_k,t_k) \in M\times [0,T) \}$, $k=1, 2, \ldots$,
is a sequence of {\bf (almost) maximum points}\index{(almost)
maximum points} if there exist positive constants $c_1$ and
$\alpha \in (0,1]$ such that
$$
|Rm(x_k,t_k)|\geq c_1 K_{\max}(t), \quad t\in
[t_k-\frac{\alpha}{K_{\max}(t_k)},t_k]
$$
for all $k$.
\end{definition}
\begin{definition}[injectivity-radius condition at almost maximum points]
  \label{def:cz-injectivity-radius-condition-at-almost-maximum-points}
  \source{cao-zhu}
  \dcref{Definition 4.3.2}
  \group{Cao--Zhu Limiting Singularity Models}
We say that the solution satisfies {\bf injectivity radius
condition}\index{injectivity radius!condition} if for any sequence
of (almost) maximum points $\{(x_k,t_k)\}$, there exists a
constant $c_2>0$ independent of $k$ such that
$$
\inj(M,x_k,g_{ij}(t_k))\geq \frac{c_2}{\sqrt{K_{\max}(t_k)}} \quad
\text{for all }\;   k.
$$
\end{definition}
\begin{definition}[singularity types for maximal Ricci flows]
  \label{def:cz-singularity-types-for-maximal-ricci-flows}
  \source{cao-zhu}
  \dcref{Definition 4.3.3}
  \group{Cao--Zhu Limiting Singularity Models}
A solution $g_{ij}(x,t)$ to the Ricci flow on the manifold $M$,
where either $M$ is compact or at each time $t$ the metric
$g_{ij}(\cdot,t)$ is complete and has bounded curvature, is called
a {\bf singularity model}\index{singularity model} if it is not
flat and of one of the following three types:

\vskip 0.1cm\noindent \textbf{Type I}\index{type I}: The solution
exists for $t\in(-\infty,\Omega)$ for some constant $\Omega$ with
$0<\Omega<+\infty$ and
$$
|Rm|\leq\Omega/(\Omega-t)
$$
everywhere with equality somewhere at $t=0$;

\vskip 0.1cm\noindent \textbf{Type II}\index{type II}: The
solution exists for $t\in(-\infty,+\infty)$  and
$$
|Rm|\leq 1 \ \ \ \ \ \ \ \ \
$$
everywhere with equality somewhere at $t=0$;

\vskip 0.1cm\noindent \textbf{Type III}\index{type III}: The
solution exists for $t\in(-A,+\infty)$ for some constant $A$ with
$0<A<+\infty$ and
$$
\ \ |Rm|\leq A/(A+t)
$$
everywhere with equality somewhere at $t=0$.
\end{definition}
\begin{theorem}[existence of dilation-limit singularity models]
  \label{thm:cz-existence-of-dilation-limit-singularity-models}
  \source{cao-zhu}
  \dcref{Theorem 4.3.4}
  \group{Cao--Zhu Limiting Singularity Models}
For any maximal solution to the Ricci flow which satisfies the
injectivity radius condition and is of Type {\rm I, II(a), (b),}
or {\rm III(a),} there exists a sequence of dilations of the
solution along (almost) maximum points which converges in the
$C_{\rm loc}^{\infty}$ topology to a singularity model of the
corresponding type.
\end{theorem}
\begin{corollary}[restricted singularity-model types under nonnegative curvature]
  \label{cor:cz-restricted-singularity-model-types-under-nonnegative-curvature}
  \source{cao-zhu}
  \dcref{Corollary 4.3.5}
  \group{Cao--Zhu Limiting Singularity Models}
For any complete maximal solution to the Ricci flow satisfying the
injectivity radius condition and with bounded and nonnegative
curvature operator on a Riemannian manifold, or on a K\"ahler
manifold with bounded and nonnegative holomorphic bisectional
curvature, there exists a sequence of dilations of the solution
along $($almost$\,)$ maximum points which converges to a singular
model.

\vskip 0.1cm \noindent For Type {\rm I} solutions:\  the limit
model exists for $t\in (-\infty,\Omega)$ with $0<\Omega<+\infty$
and has
$$
R\leq \Omega/(\Omega-t)
$$
everywhere with equality somewhere at $t=0$.

\vskip 0.1cm \noindent For Type {\rm II} solutions:\ the limit
model exists for $t\in (-\infty, +\infty)$  and has
$$
R\leq 1
$$
everywhere with equality somewhere at $t=0$.

\vskip 0.1cm \noindent For Type {\rm III} solutions:\ the limit
model exists for $t\in(-A, +\infty)$ with $0<A<+\infty$ and has
$$
R\leq A/(A+t)
$$
everywhere with equality somewhere at $t=0$.
%$$\eqno\#$$\vskip 0.3cm
\end{corollary}
\begin{theorem}[Type II models with positive Ricci curvature are steady solitons]
  \label{thm:cz-type-ii-models-with-positive-ricci-curvature-are-steady-solitons}
  \source{cao-zhu}
  \uses{thm:cz-matrix-liyauhamilton-estimate}
  \dcref{Theorem 4.3.6}
  \group{Cao--Zhu Limiting Singularity Models}
\
\begin{itemize}
\item[(i)] \ {\rm (Hamilton \cite{Ha93E})} \ Any Type {\rm II}
singularity model with nonnegative curvature operator and positive
Ricci curvature to the Ricci flow on a manifold $M$ must be a
$($steady$)$ Ricci soliton. \item[(ii)] \ {\rm (Chen-Zhu
\cite{CZ00})} \ Any Type {\rm III} singularity model with
nonnegative curvature operator and positive Ricci curvature on a
manifold $M$ must be a homothetically expanding Ricci soliton.
\end{itemize}
\end{theorem}
\begin{remark}[noncompact Type II models and steady solitons]
  \label{rem:cz-noncompact-type-ii-models-and-steady-solitons}
  \source{cao-zhu}
  \uses{thm:cz-type-ii-models-with-positive-ricci-curvature-are-steady-solitons}
  \dcref{Remark 4.3.7}
  \group{Cao--Zhu Limiting Singularity Models}
Recall from Section 1.5 that any compact steady Ricci soliton or
expanding Ricci soliton must be Einstein. If the manifold $M$ in
Theorem~\ref{thm:cz-type-ii-models-with-positive-ricci-curvature-are-steady-solitons} is noncompact and simply connected, then the steady
(or expanding) Ricci soliton must be a steady (or expanding)
gradient Ricci soliton. For example, we know that $\nabla_iV_j$ is
symmetric from (4.3.15).  Also, by the simply connectedness of $M$
there exists a function $F$ such that
$$
\nabla_i\nabla_jF=\nabla_iV_j,\qquad \text{on }\;M.
$$
So
$$
R_{ij}=\nabla_i\nabla_jF-\frac{g_{ij}}{2t},\qquad \text{on }\;M
$$
This means that $g_{ij}$ is an expanding gradient Ricci soliton.
\end{remark}
\begin{theorem}[Kähler Type II models are steady solitons]
  \label{thm:cz-kahler-type-ii-models-are-steady-solitons}
  \source{cao-zhu}
  \dcref{Theorem 4.3.8}
  \group{Cao--Zhu Limiting Singularity Models}
 \
\begin{itemize}
\item[(i)] Any Type {\rm II} singularity model on a K\"ahler
manifold with nonnegative holomorphic bisectional curvature and
positive Ricci curvature must be a steady K\"ahler-Ricci soliton.
\item[(ii)] Any Type {\rm III} singularity model on a K\"ahler
manifold with nonnegative holomorphic bisectional curvature and
positive Ricci curvature must be an expanding K\"ahler-Ricci
soliton.
\end{itemize}
\end{theorem}
\begin{theorem}[compact Type I limits are shrinking solitons]
  \label{thm:cz-compact-type-i-limits-are-shrinking-solitons}
  \source{cao-zhu}
  \dcref{Theorem 4.3.9}
  \group{Cao--Zhu Limiting Singularity Models}
Let $(M,g_{ij}(t))$ be a Type {\rm I} singularity model obtained
as a rescaling limit of a Type {\rm I} maximal solution. Suppose
$M$ is compact. Then $(M,g_{ij}(t))$ must be a gradient shrinking
Ricci soliton.
\end{theorem}
\section{Ricci Solitons}
\begin{lemma}[scalar-curvature identities for steady gradient solitons]
  \label{lem:cz-scalar-curvature-identities-for-steady-gradient-solitons}
  \source{cao-zhu}
  \dcref{Lemma 4.4.1}
  \group{Cao--Zhu Ricci Solitons}
Suppose we have a complete gradient steady Ricci soliton $g_{ij}$
with bounded curvature so that
$$
R_{ij}=\nabla_i\nabla_jF
$$
for some function $F$ on $M$. Assume the Ricci curvature is
positive and the scalar curvature $R$ attains its maximum
$R_{\max}$ at a point $x_0\in M$. Then \be
|\nabla F|^2+R=R_{\max} %\eqno(4.4.1)
\ee everywhere on $M$, and furthermore $F$ is convex and attains
its minimum at $x_0$.
\end{lemma}
\begin{theorem}[asymptotic scalar curvature of positively curved steady solitons]
  \label{thm:cz-asymptotic-scalar-curvature-of-positively-curved-steady-solitons}
  \source{cao-zhu}
  \uses{lem:cz-scalar-curvature-identities-for-steady-gradient-solitons}
  \dcref{Theorem 4.4.2}
  \group{Cao--Zhu Ricci Solitons}
For a complete noncompact steady gradient Ricci soliton with
bounded curvature and positive sectional curvature of dimension
$n\ge3$ where the scalar curvature assume its maximum at a point
$O\in M$, the asymptotic scalar curvature ratio is infinite, i.e.,
$$
A=\limsup_{s\rightarrow +\infty}Rs^2=+\infty
$$
where s is the distance to the point $O$.
\end{theorem}
\begin{theorem}[classification of two-dimensional steady solitons]
  \label{thm:cz-classification-of-two-dimensional-steady-solitons}
  \source{cao-zhu}
  \dcref{Theorem 4.4.3}
  \group{Cao--Zhu Ricci Solitons}
The only complete steady Ricci soliton on a two-dimensional
manifold with bounded curvature which assumes its maximum $1$ at
an origin is the ``cigar" soliton on the plane $\mathbb{R}^2$ with
the metric
$$
ds^2=\frac{dx^2+dy^2}{1+x^2+y^2}.
$$
\end{theorem}
